\section{The largest one-step fibres}

Let $L_0=2$, $L_1=1$, $L_{t+2}=L_{t+1}+L_t$. Equation~\eqref{eq:complement}
allows us to count the strict-lower map $F$; the targets and their fibre
sizes are unchanged for $U$. The full inverse sources for $U$ are obtained
by complementing each source constructed below.

\begin{theorem}[Sharp maximum and all equality targets]\label{thm:max}
 For every $n\ge3$,
 \[
 \max_{b\in X_n}|U^{-1}(b)|=L_{2\lfloor n/2\rfloor}.
 \]
 For even $n=2m\ge4$, equality holds exactly at the two rotations of
 $(02)^m$. For odd $n=2m+1\ge5$, it holds exactly at all rotations of
 \[
 00(20)^{m-1}2\quad\text{and}\quad011(02)^{m-1},
 \]
 giving $2n$ labelled targets. For $n=3$, the seven equality targets
 are $000$ and the rotations of $002$ and $011$.
\end{theorem}

\subsection{A complete static decoder}
Every $F$-image has a zero at a source minimum. The fibre of $0^n$
consists of the three constant sources, because zero outputs at both
ends of an edge force equality of its source heights. Otherwise, write
the target's cyclic zero runs as $Z_1,\ldots,Z_r$ and its intervening
nonempty positive runs as $w_1,\ldots,w_r$. The source is constant,
say $a_j\in\{0,1,2\}$, on $Z_j$. If $u$ is the source on $w_j$,
the two boundary inequalities are $a_j\le u_1$ and
$a_{j+1}\le u_{|u|}$, with $a_{r+1}=a_1$.

A positive target position cannot have source height zero. An interior
source one in a positive run would have two positive neighbors and
output zero. Thus all interior source heights are two. Length five
would contain three interior twos and give a zero at the middle one,
so positive runs have length at most four. For outside heights $a,b$,
the following list gives every possible positive source string $u$.
Overlapping conditions on different rows describe distinct sources.
\begin{center}\small
\begin{tabular}{lll}
\toprule
Target $w$ & Source $u$ & Outside condition\\
\midrule
$2$ & $1$; $2$ & $a=b=0$; $a,b\in\{0,1\}$\\
$1$ & $1$; $2$ & $(a,b)=(0,1),(1,0)$; $(a,b)=(0,2),(1,2),(2,0),(2,1)$\\
$11$ & $11$; $22$ & $a=b=0$; $a,b\in\{0,1\}$\\
 & $12$; $21$ & $(a,b)=(0,2)$; $(a,b)=(2,0)$\\
$12$; $21$ & $12$; $21$ & $a=0,\ b\le1$; $a\le1,\ b=0$\\
$111$ & $122$; $221$ & $a=0,\ b\le1$; $a\le1,\ b=0$\\
$121$; $1111$ & $121$; $1221$ & $a=b=0$ in either case\\
\bottomrule
\end{tabular}
\end{center}
For length one this checks heights one and two; for length two it
checks $11,12,21,22$. At length three the middle is two, $222$ fails,
and the remaining possibilities are $122,221,121$. At length four the
middle pair is $22$ and positivity there forces both endpoints to one.
This proves completeness of the list. Every unlisted target run is
impossible.

Let $K_w[a,b]$ count the listed strings, with row and column order
$0,1,2$. The eight nonzero word kernels are
\begin{gather*}
 A=K_2=\begin{pmatrix}2&1&0\\1&1&0\\0&0&0\end{pmatrix},\quad
 J_0=K_1=\begin{pmatrix}0&1&1\\1&0&1\\1&1&0\end{pmatrix},\quad
 B_0=K_{11}=\begin{pmatrix}2&1&1\\1&1&0\\1&0&0\end{pmatrix},\\
 D_0=K_{12}=\begin{pmatrix}1&1&0\\0&0&0\\0&0&0\end{pmatrix},\quad
 K_{21}=D_0^{\mathsf T},\quad
 C_0=K_{111}=\begin{pmatrix}2&1&0\\1&0&0\\0&0&0\end{pmatrix},\\
 K_{121}=K_{1111}=E_0=\begin{pmatrix}1&0&0\\0&0&0\\0&0&0\end{pmatrix}.
\end{gather*}
The subscripts distinguish $J_0$ from input complement $J$, and $B_0$
from the generating-function matrix $B(z)$. Choose $a_1,\ldots,a_r$,
fill the actual zero-run coordinates with those heights, and independently
choose one listed source string for each positive run. The boundary
inequalities enforce the required zero outputs. Conversely every source
recovers exactly these choices. This labelled bijection gives
\begin{equation}\label{eq:fibretrace}
 |F^{-1}(b)|=\sum_{a_1,\ldots,a_r=0}^2
 \prod_{j=1}^r K_{w_j}[a_j,a_{j+1}]
 =\tr(K_{w_1}\cdots K_{w_r}).
\end{equation}
It includes $r=1$, when both outside heights coincide. A different
start cyclically permutes factors, without multiplying by rotations.
Zero-run lengths affect reconstructed coordinates but not the kernels.

This decoder is static background. More generally, the source edge signs
are recoverable; contracting their zero edges and orienting every
remaining edge from lower to higher gives strict order-map sets on the
resulting posets. The fibre is the disjoint union over compatible sign
words, with no lost labels. The issue in Theorem~\ref{thm:max} is to
compare the sizes of these \emph{whole} unions across rank targets.

\subsection{Uniform comparison of mixed kernels}
Set $\lambda=(3+\sqrt5)/2$. We use the standard Schatten Hölder bound
\begin{equation}\label{eq:holder}
 |\tr(M_1\cdots M_r)|\le\prod_{j=1}^r\snorm{M_j}{r}
 \qquad(r\ge2)
\end{equation}
for real square matrices of equal size, with no commutation or
invertibility hypothesis. Here $\snorm{M}{p}$ is the $\ell^p$ norm
of its singular values. To specify the source and exponents precisely,
the two-factor unitarily invariant norm inequality in
\citet[Example 6.18 and Theorem 6.32]{tropp2022matrix} yields
$\snorm{VW}{t}\le\snorm{V}{pt}\snorm{W}{qt}$ for $t\ge1$,
$p,q>1$ conjugate. In the induction on the first $s$ factors use
$t=r/s$, $p=s/(s-1)$, $q=s$, for $2\le s\le r$.
Finally $|\tr P|\le\snorm{P}{1}$ follows from an SVD, since the
diagonal entries of the intervening orthogonal matrix have absolute
value at most one. This proves~\eqref{eq:holder} from that standard
inequality with valid exponents throughout.

All kernels have nonnegative entries, and
$D_0,D_0^{\mathsf T},C_0,E_0\le A$ entrywise. Replacing any of the
five corresponding target-word kernels by $A$ cannot decrease the trace
in~\eqref{eq:fibretrace}: expand it as a sum of products of entries.
No positive-semidefinite order claim is involved. In a resulting length-$r$
product, let $k$ count $B_0$ factors and $j$ count $J_0$ factors.
The positive eigenvalues of $A$ are $\lambda,\lambda^{-1}$; write
\[
 a_r=\snorm{A}{r}=(\lambda^r+\lambda^{-r})^{1/r}.
\]
The singular values of $J_0$ are $2,1,1$, and
$\snorm{B_0}{2}=3$. Monotonicity of vector norms gives
\begin{equation}\label{eq:normbounds}
 \snorm{J_0}{r}\le\sqrt6<\lambda<a_r,\qquad
 \snorm{B_0}{r}\le3\qquad(r\ge2).
\end{equation}
For $r\ge3$ we also have $\snorm{J_0}{r}\le\sqrt[3]{10}$ and
$3\sqrt[3]{10}<\lambda^2\le a_r^2$. For instance
$\lambda>13/5$ and $169^3>270\cdot25^3$ prove the strict cubic
inequality; $\lambda^2=(7+3\sqrt5)/2>6$ proves the first strict
comparison in~\eqref{eq:normbounds}.

If $k=0$, Hölder bounds the trace by
$a_r^r=\lambda^r+\lambda^{-r}$, strictly when $j>0$.
If $k=1$ and $j=0$, cyclicity and the leading two-by-two block give
\begin{equation}\label{eq:oneb}
 \tr(B_0A^{r-1})=\tr A^r=\lambda^r+\lambda^{-r},
\end{equation}
because $r-1\ge1$ removes the third-coordinate contribution. If
$k=1$ and $j>0$, the case $r=2$ gives
$\tr(B_0J_0)=4<7=\tr A^2$. For $r\ge3$, combine one $B_0$ and
one $J_0$ in the \emph{scalar product of norm bounds} using
$3\sqrt[3]{10}<a_r^2$; every additional $J_0$ has smaller norm than
$A$. The resulting bound is again strict below $a_r^r$.

For $k\ge2$, Hölder instead gives
\begin{equation}\label{eq:manyb}
 \tr(M_1\cdots M_r)\le a_r^{r-k}3^k
 <\frac{10}{9}\lambda^{r-k}3^k
 <\lambda^{r+\lfloor k/2\rfloor}.
\end{equation}
Indeed $(a_r/\lambda)^r=1+\lambda^{-2r}<10/9$, since
$r\ge2$ and $\lambda^4>9$; this also handles $r=k$.
The last inequality is equivalent to
$(10/9)3^k<\lambda^{k+\lfloor k/2\rfloor}$. For $k=2,3$ it is
$10<\lambda^3$ and $30<\lambda^4$; increasing $k$ by two multiplies
the two sides by $9$ and $\lambda^3>9$, respectively. This proves
every mixed-product bound without commuting the factors.

\subsection{Length budget and complete equality analysis}
Each positive run and a following nonempty zero run use at least two
target positions, and each $B_0=K_{11}$ costs one extra. Consequently
\begin{equation}\label{eq:budget}
 n\ge2r+k.
\end{equation}
The even Lucas subsequence satisfies
$L_{2s}=\lambda^s+\lambda^{-s}$, since both sides start at $2,3$
and satisfy $V_{s+2}=3V_{s+1}-V_s$; it is strictly increasing for
$s\ge1$. If $r\ge2$ and $k\le1$, the preceding bounds give a
trace at most $L_{2r}\le L_{2\lfloor n/2\rfloor}$. For $k\ge2$,
\eqref{eq:manyb} and~\eqref{eq:budget} give a strict bound below
$L_{2\lfloor n/2\rfloor}$.

Targets without zeros have no source. The all-zero target has fibre
three. For $r=1$, the traces of $J_0,A,B_0,D_0,D_0^{\mathsf T},C_0,E_0$
are $0,3,3,1,1,2,1$; all other kernels vanish. These are below
$L_4=7$ for $n\ge4$. At $n=3$, at most one positive run exists,
so precisely $000$ and rotations of $002,011$ attain three. This
proves the bound and all small-length cases.

For equality with $n\ge4$, we must have $r\ge2$, $k\le1$, and
$q:=n-2r\in\{0,1\}$: if $q\ge2$, then
$L_{2r}<L_{2\lfloor n/2\rfloor}$. If $q=0$, all positive and zero
runs have length one. A positive run one gives $J_0$ and strictness,
so equality holds exactly for alternating $02$ targets, which attain
$\tr A^r=L_{2r}$.

If $q=1$, exactly one run has one extra position. If this is a zero
run, all positive runs must again be two and the targets are the
rotations of $00(20)^{r-1}2$. If it is a positive run, its nonzero
kernel is $B_0,D_0$, or $D_0^{\mathsf T}$. For $B_0$, any $J_0$
factor gives strictness, and all remaining factors $A$ give equality
by~\eqref{eq:oneb}. This is the family $011(02)^{r-1}$.
For $D_0$ or $D_0^{\mathsf T}$, a $J_0$ factor again makes the
entrywise-replaced bound strict. Without $J_0$, the leading
two-by-two block of $A^{r-1}$ is strictly positive, while $A-D_0$
and $A-D_0^{\mathsf T}$ are nonnegative and nonzero in that block.
The trace expansion yields a strict inequality below $\tr A^r$.
No longer run fits the single extra position. This exhausts all
equality cases and proves Theorem~\ref{thm:max}. The two odd families
each have a unique doubled run, forbidding a nontrivial rotational
stabilizer; the presence of ones distinguishes them. Each therefore
gives $n$ labelled targets.

For completeness, the attainer's classical counting mechanism is
explicit. For an alternating lower-rank target, source valleys have
heights in $\{0,1\}$, peaks in $\{1,2\}$, and every valley is
strictly lower than its adjacent peaks. Mark exactly the sites of
source height one. These marks are an independent set of the labelled
cycle, and the inverse construction puts zero on unmarked valleys,
two on unmarked peaks, and one on marked sites. Its count is $L_{2m}$.
Doubling a valley source height gives the odd $00$ family. Doubling
a peak gives the odd $11$ family: its boundary valleys are in
$\{0,1\}$, so the local source choices are precisely $11$ or $22$.
Deleting the repeated position reverses either construction. These
full set bijections explain attainment, but the mixed-target
comparison above is needed to prove global optimality.
